\documentclass{situationist-preprint}

\title{Bentham's Objection Reconsidered}
\msid{POL-18-b}
\msfamily{POL-18}
\msdiscipline{Political Philosophy}
\mstype{Companion manuscript}
\mskeywords{rights; Bentham; nonsense upon stilts; paraphrasis; fictitious entities; Hohfeld}

\begin{document}
\maketitle

\begin{abstract}
POL-18-a answers Bentham's attack on natural rights by rejecting the command theory on which it rests. A stronger version of the objection does not depend on that theory. Bentham held that \enquote{right} names a fictitious entity, whose terms have meaning only through paraphrase into sentences about real entities, and that the only available paraphrase of rights-sentences is into sentences about legal obligations and sanctions. We argue that the method of paraphrasis is sound and that rights-sentences survive it: they paraphrase into sentences about reasons, owed, waivable and exclusionary, and Lemma~3.1 of POL-18-a shows that reasons so structured are what rights are. The stronger objection therefore fails for the same reason as the weaker one.
\end{abstract}

\section{The stronger objection}

In \emph{Anarchical Fallacies}, his commentary on the French Declaration of Rights, Bentham argued that rights are the fruit of law and of law alone, and that natural rights are nonsense \citep{Bentham1843}. POL-18-a reads the argument through the command theory, on which a right exists only where a sovereign commands and a sanction is attached, and replies with Hart's demonstration that the command theory fails even for legal rights \citep{Hart1982}.

Bentham's semantic doctrine supports a version of the objection that does not need the command theory. Many nouns, on his view, name fictitious entities: they do not denote anything, and they are not for that reason meaningless, but they have meaning only through \emph{paraphrasis}, the translation of whole sentences in which they occur into sentences about real entities. \enquote{Obligation} and \enquote{right} are such terms. A sentence ascribing a right is meaningful if it paraphrases into a sentence about what is real; and Bentham held that the real entities to which rights-talk reduces are pains, pleasures and the sanctions attached to them by a lawgiver \citep{Hart1982}. On this version the question is not whether a command is needed but whether rights-sentences have a paraphrase that makes them true without a law.

\section{The paraphrase}

We accept the method. The claim of POL-18-a is not that \enquote{right} denotes a simple object but that rights-sentences are true. By Definition~2.1 of POL-18-a, \enquote{$A$ has a claim-right against $B$ that $B$ not $\phi$} paraphrases as: there is a reason for $B$ not to $\phi$, owed to $A$, waivable by $A$, and excluding from $B$'s deliberation certain reasons for $\phi$-ing. This is a paraphrasis in Bentham's sense: the whole sentence is translated into one whose terms are reasons and persons, and the right is not posited as an additional thing. It also matches the correlativity that \citet{Hohfeld1913} found in legal relations, of a claim in one person with a duty in another.

\begin{proposition}\label{p:survive}
Rights-sentences survive Bentham's method: they have paraphrases in terms of reasons, which are true in some cases where no law exists.
\end{proposition}

\begin{proof}
By Assumption~2.4 of POL-18-a reasons are real, so a paraphrase into sentences about reasons is a paraphrase into sentences about what is real. By Lemma~3.1 of POL-18-a a claim-right is such a structure of reasons, and by Lemma~3.2 some such structures obtain independently of any law.
\end{proof}

The disagreement with Bentham is thus located exactly. It is not about whether \enquote{right} is a fictitious term; it is about which terms are fit to appear on the right-hand side of a paraphrase. Bentham admitted only pleasures, pains and sanctions. Lemma~3.3 of POL-18-a shows that this basis is too poor even for the phenomena Bentham accepted: when a right-holder waives her right, what $B$ may do changes and no pleasure, pain or sanction need change.

\begin{objection}
Bentham's restriction of the basis to pleasures and pains is what made his method a tool of criticism. A paraphrase in terms of reasons lets every declaration of rights pass.
\end{objection}

\begin{reply}
It does not. A rights-sentence is true only if the reasons it describes exist, and whether they do is a substantive question. Bentham's practical complaint, that declarations of indefeasible rights are anarchical because they deny every law that limits them, retains its force against declarations of that form; POL-18-a answers it by the defeasibility of rights, not by the method of paraphrase \citep{Waldron1987}.
\end{reply}

\section{Conclusion}

Reconsidered in its strongest form, Bentham's objection is a demand for paraphrase. Rights meet it: rights-sentences translate into sentences about reasons, and some of those sentences are true without any law. Rights are real normative constraints, as the conclusion of POL-18 holds, and they are so on Bentham's own conditions for meaningful discourse once reasons are admitted to the basis.

\begin{thebibliography}{9}\small
\bibitem[Bentham(1843)]{Bentham1843} Bentham, J. (1843). Anarchical fallacies. In J.~Bowring (ed.), \emph{The Works of Jeremy Bentham}, vol.~2. William Tait.
\bibitem[Hart(1982)]{Hart1982} Hart, H. L. A. (1982). \emph{Essays on Bentham: Studies in Jurisprudence and Political Theory}. Clarendon Press.
\bibitem[Hohfeld(1913)]{Hohfeld1913} Hohfeld, W. N. (1913). Some fundamental legal conceptions as applied in judicial reasoning. \emph{Yale Law Journal}, 23(1), 16--59.
\bibitem[Waldron(1987)]{Waldron1987} Waldron, J. (ed.) (1987). \emph{\enquote{Nonsense upon Stilts}: Bentham, Burke and Marx on the Rights of Man}. Methuen.
\end{thebibliography}

\end{document}
